%%%%%%%%%%%%%%%%%%%%%%%%%%%%%%%%%%%%%%%
% Trabalho de Sistemas Distribu�dos
%%%%%%%%%%%%%%%%%%%%%%%%%%%%%%%%%%%%%%%

\documentclass[11pt,brazil, a4paper]{article}

%%%%%%%%%%%%%%%%%%%%%%%%%%%%%%%%%%%%%%%
% Preambulo - Configura��o de p�gina
%%%%%%%%%%%%%%%%%%%%%%%%%%%%%%%%%%%%%%%
\usepackage{color}
\usepackage[portuges]{babel}
\usepackage[T1]{fontenc}
\usepackage[ansinew]{inputenc}
\usepackage{array}
\usepackage{colortbl}
\usepackage{epsfig}
\usepackage{multirow}
\usepackage{hyperref}
\usepackage{amsmath}

\usepackage{tikz}
\usetikzlibrary{3d}
\usetikzlibrary{calc,decorations.pathmorphing} % noisy shapes
\usetikzlibrary{fit}					% fitting shapes to coordinates
\usetikzlibrary{backgrounds}	% drawing the background after the foreground
\usetikzlibrary{shadows}
\usetikzlibrary{mindmap}
\usetikzlibrary{positioning,shapes,shadows,arrows}
\usepackage{pgfplots}
\pgfplotsset{compat=1.8}
\usetikzlibrary{arrows,shapes}

\pagestyle{empty}

% Define margem no topo como sendo 1,7 cm
\setlength{\topmargin}{0pt}
\setlength{\voffset}{-1.8cm}

%define margem inferior como 1,5 cm
\setlength{\textheight}{25,6cm}

%define margens laterais como sendo 1,5 cada
\setlength{\oddsidemargin}{0cm}
\setlength{\hoffset}{-1.2cm}
\setlength{\textwidth}{18cm}

%Define campo com sombreamento cinza
\newcolumntype{G}{>{\columncolor[gray]{0.8}[\tabcolsep]\tiny \bf \sf}p{13pt}<{\centering}}
\newcolumntype{T}{>{\tiny \sf}p{13pt}<{\centering}}


\usepackage{xcolor}
\definecolor{verde}{rgb}{0,0.5,0}

\usepackage{listings}
\lstset{
	language=c,
	basicstyle=\ttfamily\small, 
	keywordstyle=\color{blue}, 
	stringstyle=\color{verde}, 
	commentstyle=\color{red}, 
	extendedchars=true, 
	showspaces=false, 
	showstringspaces=false, 
	numbers=left,
	numberstyle=\tiny,
	breaklines=true, 
	backgroundcolor=\color{green!10},
	breakautoindent=true, 
	captionpos=b,
	xleftmargin=0pt,
}


%%%%%%%%%%%%%%%%%%%%%%%%%%%%%%%%%%%%%%%
% Inicio do documento
%%%%%%%%%%%%%%%%%%%%%%%%%%%%%%%%%%%%%%%
\begin{document}
	
	\selectlanguage{brazil}
	
	\hspace{-0.8cm} \begin{tabular*}{18cm}{p{2cm}p{12cm}@{\extracolsep{\fill}}}
		\multirow{2}{2cm}{\epsfig{file=puclogo_small_bw,width=1.5cm}}
		&  \vspace{1mm} \usefont{T1}{cmss}{bx}{n} {\centering PONTIF�CIA UNIVERSIDADE CAT�LICA DE MINAS GERAIS }\\
		&  \usefont{T1}{cmss}{x}{n} Instituto de Ci�ncias Exatas e Inform�tica \vspace{3mm}  \\
	\end{tabular*}
	
	
	%%%%%%%%%%%%%%%%%%%%%%%%%%%%%%%%%%%%%%%
	% Define caracteristicas da disciplina
	%%%%%%%%%%%%%%%%%%%%%%%%%%%%%%%%%%%%%%%
	\hspace{-0.8cm} \begin{tabular*}{18.5cm}{|p{7cm}@{\extracolsep{\fill}}|p{6cm}|p{1.7cm}|}
		\hline
		\sf Disciplina & \sf Curso  & \sf Per�odo \\
		
		\multicolumn{1}{|r|}{Algoritmos e Estruturas de Dados II} &
		\multicolumn{1}{c|}{Engenharia de Computa��o} &
		\multicolumn{1}{c|}{2$^\circ$} \\ \hline \hline
		\multicolumn{3}{|l|}{\sf Professor} \\
		\multicolumn{3}{|l|}{\hspace{1cm}Kleber Jacques F. de Souza (klebersouza@pucminas.br)} \\ \hline
	\end{tabular*}
	
	
	\vspace{0.1cm}
	
	
	\begin{center}
		{\Large \bf Lista de Exerc�cios - Recursividade} \\
		%{\normalsize \bf em dupla}
	\end{center}
	

Desenvolva algoritmos recursivos para os seguintes problemas:

\begin{enumerate}

\item Multiplica��o de dois n�meros naturais, atrav�s de somas sucessivas (Ex.: 6 * 4 = 4+4+4+4+4+4).

\begin{lstlisting}
int multiplicacao (int a, int b){
	if(b == 0) return 0;
	if(b < 0){
		a *= -1;
		b *= -1;
	}	
	return (a + multiplicacao (a, b-1));
}
\end{lstlisting}

\item C�lculo de 1 + $\frac{1}{2}$ + $\frac{1}{3}$ + $\frac{1}{4}$ + ... + $\frac{1}{N}$.

\begin{lstlisting}
float f(int n){
	if(n == 1) return 1;
	else return 1.0/n + f(n-1);
}
\end{lstlisting}

\item C�lculo de $\frac{2}{4}$ + $\frac{5}{5}$ + $\frac{10}{6}$ + $\frac{17}{7}$ + $\frac{26}{8}$ + ... + $\frac{N^2+1}{N+3}$.

\begin{lstlisting}
float f(int n){
	if(n == 1) return (2.0/4.0);
	else return ((n*n)+1)/(n+3) + f(n-1);
}
\end{lstlisting}


\item Invers�o de uma string.

\begin{lstlisting}

void inverte (char palavra[], int n){
	
	printf("%c",palavra[n]);
	
	if(n == 0) return;
	else inverte(palavra, n-1);
}

\end{lstlisting}

\item  Gerador da sequ�ncia dada por:
\begin{itemize}
\item  F(1) = 1
\item  F(2) = 2
\item  F(n) = 2 * F(n - 1) + 3 * F(n - 2).
\end{itemize}

\begin{lstlisting}
int sequencia(int n){
	if(n <= 2) return n;
	else return 2*sequencia(n-1) + 3 *sequencia(n-2);
}
\end{lstlisting}

\item A partir de um vetor de n�meros inteiros, calcule a soma e o produto dos elementos
do vetor.


\item  Gerador de m�ximo divisor comum (mdc):
\begin{itemize}
\item  mdc(x, y) = y, se x $\geq$ y e x mod y = 0
\item  mdc(x, y) = mdc(y, x), se x < y
\item  mdc(x, y) = mdc(y, x mod y), caso contr�rio
\end{itemize}

\begin{lstlisting}
int mdc (int a, int b) {
	if (b == 0) return a;
	else return mdc(b, a%b);
}
\end{lstlisting}

\item Verifique se uma palavra � pal�ndromo (Ex. aba, abcba, xyzzyx).

\begin{lstlisting}
bool palindromo(char palavra[], int i, int f){

	if(i>=f)
	return true;
	
	if(palavra[i]==palavra[f])
	palindromo(palavra, i+1, f-1);
	else
	return false;
}
\end{lstlisting}

\item Dado um n�mero n, gere todas as poss�veis combina��es com as n primeiras letras do alfabeto.\\ Ex.: n = 3. Resposta: ABC, ACB, BAC, BCA, CAB, CBA.

\begin{lstlisting}
void troca(char vetor[], int i, int j)	{
	char aux = vetor[i];
	vetor[i] = vetor[j];
	vetor[j] = aux;
}

void permuta(char vetor[], int inf, int sup)	{
	if(inf == sup)	{
		for(int i = 0; i <= sup; i++)
			printf("%c ", vetor[i]);
		printf("\n");
	}
	else{
		for(int i = inf; i <= sup; i++)	{
			troca(vetor, inf, i);
			permuta(vetor, inf + 1, sup);
			troca(vetor, inf, i);
		}
	}
}

int main(int argc, char *argv[])	{
	char alfabeto[]="ABCDEFGHIJKLMNOPQRSTUVWXYZ";
	int N;
	
	cout << "Digite o valor N: ";
	cin >> N;
	
	permuta(alfabeto, 0, N - 1);
	
	return 0;
}
\end{lstlisting}

\item Gere todas as poss�veis combina��es para um jogo da MegaSena com 6 dezenas.


\begin{lstlisting}
#include <stdio.h>
#include <stdlib.h>

void combinations (int v[], int start, int n, int k, int maxk) {
	int     i;
	
	/* k here counts through positions in the maxk-element v.
	* if k > maxk, then the v is complete and we can use it.
	*/
	if (k > maxk) {
		/* insert code here to use combinations as you please */
		for (i=1; i<=maxk; i++) printf ("%i ", v[i]);
			printf ("\n");
		return;
	}
	
	/* for this k'th element of the v, try all start..n
	* elements in that position
	*/
	for (i=start; i<=n; i++) {
		
		v[k] = i;
		
		/* recursively generate combinations of integers
		* from i+1..n
		*/
		combinations (v, i+1, n, k+1, maxk);
	}
}
	
int main (int argc, char *argv[]) {
	int     v[100], n, k;
	
	n = 60;
	k = 6;
	
	/* generate all combinations of n elements taken
	* k at a time, starting with combinations containing 1
	* in the first position.
	*/
	combinations (v, 1, n, 1, k);
	exit (0);
}

\end{lstlisting}

\item Implemente uma fun��o recursiva que, dados dois n�meros inteiros $x$ e $n$, calcula o valor de $x^n$. 



\begin{lstlisting}
int potencia (int base, int expoente){
	if(expoente == 0)
		return 1;
	else
		return (base * potencia(base, expoente-1));
}
\end{lstlisting}

\end{enumerate}




%===================================================================================================

\end{document}